\section{Experimentación y Resultados}
\label{sec:evidencia}

Se ejecutó un experimento computacional comparativo para observar cómo cambios
en estabilidad y eficiencia modifican el desempeño durante una operación
industrial. Los cuatro casos utilizan el protocolo definido en la
Sección \ref{sec:metodologia} y mecanismos consistentes con la literatura
\citep{siddiqui_defying_2017,hou_tradeoff_2023}.

\subsection{Resultados por Aplicación}
\label{subsec:resultados_aplicacion}

La Figura \ref{fig:produccion_casos} muestra el resultado productivo bajo la
condición base: \(60\,^{\circ}\mathrm{C}\), 40 mM de sustrato inicial y ocho
horas de operación.

\begin{figure}[H]
\centering

\begin{tabular}{|c|c|}
\hline

\parbox{0.46\textwidth}{
    \centering
    \small Caso 1: Procesamiento de Alimentos\\[0.12cm]
    \includegraphics[width=0.43\textwidth]{13.png}
}
&
\parbox{0.46\textwidth}{
    \centering
    \small Caso 2: Bagazo de Caña\\[0.12cm]
    \includegraphics[width=0.43\textwidth]{14.png}
}
\\
\hline

\parbox{0.46\textwidth}{
    \centering
    \small Caso 3: Procesamiento de Biomasa\\[0.12cm]
    \includegraphics[width=0.43\textwidth]{15.png}
}
&
\parbox{0.46\textwidth}{
    \centering
    \small Caso 4: Síntesis de Glucósidos\\[0.12cm]
    \includegraphics[width=0.43\textwidth]{16.png}
}
\\
\hline

\end{tabular}

\figcaptionlink{anx:fig_produccion}{Producción Acumulada Bajo la Condición Base}{fig:produccion_casos}


\end{figure}

\subsubsection{Caso 1: Procesamiento de Alimentos}

La proteasa ENG conserva actividad durante más tiempo que WT y alcanza
\(4.58\times\) la producción acumulada bajo la condición base. La principal
ganancia es extender la vida operativa sin sacrificar capacidad de conversión
\citep{yuan_apre2709_2025}. Las curvas completas de estabilidad se presentan en
el \anexolink{anx:estabilidad}{Estabilidad Térmica}.

\subsubsection{Caso 2: Bagazo de Caña}

La beta-glucosidasa ENG combina mayor estabilidad con mayor tolerancia a glucosa.
El producto acumulado alcanza \(8.96\times\) el valor de WT. En este caso, la
ventaja principal es mantener la reacción activa cuando el producto comienza a
inhibirla \citep{cao_ks5a7_2018}. La evidencia de estabilidad y cinética se
presenta en los \anexolink{anx:estabilidad}{Estabilidad Térmica} y \anexolink{anx:cinetica}{Cinética Frente a Concentración de Sustrato}.

\subsubsection{Caso 3: Procesamiento de Biomasa}

La xilanasa ENG desplaza su respuesta hacia temperaturas más elevadas y alcanza
\(5.41\times\) la producción de WT bajo la condición base
\citep{zhang_xylanase_2010}. La comparación térmica completa se presenta en el
\anexolink{anx:temperatura}{Respuesta Térmica}.

\subsubsection{Caso 4: Síntesis de Glucósidos}

La variante ENG combina estabilidad, mayor eficiencia y mejor selectividad. El
efecto conjunto genera el mayor aumento bajo condición base:
\(15.29\times\) la producción de WT \citep{yin_bgl1d_2019}. La cinética y la
selectividad se presentan en los \anexolink{anx:cinetica}{Cinética Frente a Concentración de Sustrato} y
\anexolink{anx:selectividad}{Selectividad en Síntesis de Glucósidos}.

\subsection{Resumen Cuantitativo}
\label{subsec:resumen_cuantitativo}

El Cuadro \ref{tab:resumen_resultados} concentra los resultados necesarios para
la decisión posterior.

\begin{table}[H]
\centering
\tabcaptionlink{anx:tab_resumen}{Resumen Cuantitativo Bajo la Condición Base}{tab:resumen_resultados}



\renewcommand{\arraystretch}{1.35}
\setlength{\tabcolsep}{3.5pt}

\begin{tabularx}{\textwidth}{
|p{2.5cm}|
>{\centering\arraybackslash}p{2.0cm}|
>{\centering\arraybackslash}p{1.45cm}|
>{\centering\arraybackslash}p{1.45cm}|
>{\centering\arraybackslash}p{1.45cm}|
>{\centering\arraybackslash}p{2.0cm}|
X|
}
\hline
Aplicación &
\(t_{50}\) WT \(\rightarrow\) ENG &
\(I_S\) &
\(I_C\) &
\(I_P\) &
\(T_{\mathrm{opt}}\) WT \(\rightarrow\) ENG &
Lectura \\ \hline

Alimentos &
1.11 \(\rightarrow\) 3.53 h &
3.37 &
1.54 &
4.58 &
60 \(\rightarrow\) 65 \(^{\circ}\)C &
La vida operativa impulsa el resultado. \\ \hline

Bagazo &
0.75 \(\rightarrow\) 7.02 h &
6.92 &
1.48 &
8.96 &
55 \(\rightarrow\) 65 \(^{\circ}\)C &
Estabilidad y tolerancia al producto. \\ \hline

Biomasa &
0.99 \(\rightarrow\) 6.13 h &
4.86 &
1.58 &
5.41 &
65 \(\rightarrow\) 75 \(^{\circ}\)C &
Mayor ventana térmica útil. \\ \hline

Glucósidos &
1.49 \(\rightarrow\) 5.43 h &
3.19 &
2.56 &
15.29 &
55 \(\rightarrow\) 60 \(^{\circ}\)C &
Eficiencia y selectividad amplifican la mejora. \\ \hline

\end{tabularx}
\end{table}

Los índices \(I_S\), \(I_C\) e \(I_P\) son mayores que 1 en los cuatro casos.
La mejora conjunta existe, pero su efecto sobre producción es distinto según el
mecanismo dominante.

\subsection{Comparación Gráfica}
\label{subsec:comparacion_grafica}

\begin{figure}[H]
\centering

\begin{tabular}{|c|c|}
\hline

\parbox{0.46\textwidth}{
    \centering
    \small (a) Impacto sobre Producción\\[0.12cm]
    \includegraphics[width=0.43\textwidth]{20.png}
}
&
\parbox{0.46\textwidth}{
    \centering
    \small (b) Matriz Estabilidad--Eficiencia\\[0.12cm]
    \includegraphics[width=0.43\textwidth]{21.png}
}
\\
\hline

\end{tabular}

\figcaptionlink{anx:fig_sintesis}{Síntesis Comparativa de los Cuatro Casos}{fig:resumen_experimento}


\end{figure}

La comparación individual de estabilidad y eficiencia se encuentra en el
\anexolink{anx:fig_sintesis}{Síntesis Comparativa de los Cuatro Casos}. No se repiten aquí esas figuras para mantener el cuerpo
principal orientado a resultados.

\subsection{Comportamientos que Limitan la Mejora}
\label{subsec:fenomenos}

La Figura \ref{fig:fenomenos} muestra por qué una mayor estabilidad o eficiencia
no se traduce siempre de forma proporcional en producto.

\begin{figure}[H]
\centering

\begin{tabular}{|c|c|}
\hline

\parbox{0.46\textwidth}{
    \centering
    \small (a) Proteasa: Riesgo de Inactivación\\[0.12cm]
    \includegraphics[width=0.43\textwidth]{22.png}
}
&
\parbox{0.46\textwidth}{
    \centering
    \small (b) Bagazo: Inhibición por Producto\\[0.12cm]
    \includegraphics[width=0.43\textwidth]{23.png}
}
\\
\hline

\parbox{0.46\textwidth}{
    \centering
    \small (c) Biomasa: Activación Térmica Transitoria\\[0.12cm]
    \includegraphics[width=0.43\textwidth]{24.png}
}
&
\parbox{0.46\textwidth}{
    \centering
    \small (d) Glucósidos: Inhibición por Sustrato\\[0.12cm]
    \includegraphics[width=0.43\textwidth]{25.png}
}
\\
\hline

\end{tabular}

\figcaptionlink{anx:fig_restricciones}{Restricciones No Lineales Observadas en los Cuatro Casos}{fig:fenomenos}


\end{figure}

\noindent\small
Evidencia complementaria:
\anexolink{anx:fig_restricciones}{Restricciones No Lineales Observadas en los Cuatro Casos} y
\anexolink{anx:selectividad}{Selectividad en Síntesis de Glucósidos}.
\normalsize

Los mecanismos son distintos: inactivación en alimentos, inhibición por producto
en bagazo, respuesta térmica en biomasa e inhibición por exceso de sustrato en
glucósidos. En este último caso, la selectividad introduce además una segunda
condición de decisión.

\subsection{Resultado General}
\label{subsec:resultado_general}

Los resultados responden la pregunta central: mejorar estabilidad no obligó a
sacrificar eficiencia en ninguno de los cuatro casos. Sin embargo, la mejora
molecular solo genera su máximo valor cuando se opera dentro de condiciones que
eviten el nuevo cuello de botella.

La \hyperref[sec:discusion]{Sección 5} transforma estos resultados en decisiones
concretas de operación y prioridad de explotación.
